% Página 17 do PDF de Gabriel
\begin{frame}[plain,noframenumbering,label=original-17]
  \begin{tikzpicture}[remember picture,overlay,x=1bp,y=-1bp]
  \begin{scope}[shift={(current page.north west)}]
  \Fundo{17}
  \Texto{44.456}{41.894}{90.657}{8.518}{b}{34E0C8}{ANÁLISE DESCRITIVA}
  \Texto{43.039}{116.125}{63.147}{56.787}{b}{FFB014}{04}
  \Texto{43.039}{155.148}{209.804}{24.608}{b}{FFFFFF}{Análise descritiva}
  \Texto{43.039}{234.662}{46.285}{7.572}{b}{34E0C8}{EMASC 2026}
  \Texto{91.429}{234.662}{2.65}{7.572}{b}{FFB014}{\textbullet{}}
  \Texto{96.176}{234.662}{21.45}{7.572}{b}{34E0C8}{LAME}
  \Link{42.042}{135.586}{211.795}{20.866}{original-17}
  \end{scope}
  \end{tikzpicture}
\end{frame}


% Página 18 do PDF de Gabriel
\begin{frame}[plain,noframenumbering,label=original-18]
  \begin{tikzpicture}[remember picture,overlay,x=1bp,y=-1bp]
  \begin{scope}[shift={(current page.north west)}]
  \Fundo{18}
  \Texto{40.175}{29.916}{253.823}{17.983}{b}{2C2084}{Medidas de tendência central}
  \Texto{35.433}{63.506}{295.975}{8.518}{i}{737696}{Uma lista de 15.000 valores não se comunica; um número representativo, sim.}
  \Texto{34.223}{100.879}{26.822}{9.465}{b}{CC8D10}{Média}
  \Texto{34.223}{115.823}{88.001}{9.465}{r}{000000}{O ponto de equilíbrio}
  \Texto{34.223}{127.778}{50.645}{9.465}{r}{000000}{dos valores.}
  \Texto{34.223}{145.711}{32.198}{11.357}{b}{000000}{4,57 h}
  \Texto{177.266}{100.642}{37.867}{9.465}{b}{027499}{Mediana}
  \Texto{177.266}{115.586}{89.089}{9.465}{r}{000000}{O valor central da fila}
  \Texto{177.266}{127.542}{42.619}{9.465}{r}{000000}{ordenada.}
  \Texto{177.266}{145.474}{16.411}{11.357}{b}{000000}{4 h}
  \Texto{320.309}{100.439}{24.712}{9.465}{b}{007078}{Moda}
  \Texto{320.309}{115.383}{83.922}{9.465}{r}{000000}{O valor que mais se}
  \Texto{320.309}{127.338}{29.321}{9.465}{r}{000000}{repete.}
  \Texto{320.309}{145.271}{16.411}{11.357}{b}{000000}{4 h}
  \Texto{49.227}{177.332}{196.246}{8.518}{r}{000000}{Horas diárias em redes sociais de sete estudantes:}
  \Texto{249.274}{177.332}{69.899}{8.518}{b}{000000}{2, 3, 4, 4, 5, 6 e 8}
  \Texto{319.172}{177.332}{87.139}{8.518}{r}{000000}{. Com poucos valores,}
  \Texto{49.227}{188.291}{357.083}{8.518}{r}{000000}{cada cálculo pode ser conferido manualmente; o procedimento é o mesmo para 15.000 regis-}
  \Texto{49.227}{199.25}{16.44}{8.518}{r}{000000}{tros.}
  \Texto{37.21}{247.463}{40.499}{6.625}{b}{737696}{EMASC 2026}
  \Texto{80.653}{247.463}{2.319}{6.625}{r}{FFB014}{\textbullet{}}
  \Texto{85.913}{247.463}{117.821}{6.625}{r}{737696}{Liga Acadêmica Matemática \& Empresa}
  \Texto{206.682}{247.463}{2.319}{6.625}{r}{FFB014}{\textbullet{}}
  \Texto{211.943}{247.463}{13.25}{6.625}{r}{737696}{UFS}
  \Texto{397.549}{248.132}{7.367}{6.625}{b}{2C2084}{13}
  \Texto{406.02}{248.132}{10.315}{6.625}{r}{737696}{/ 40}
  \end{scope}
  \end{tikzpicture}
\end{frame}


% Página 19 do PDF de Gabriel
\begin{frame}[plain,noframenumbering,label=original-19]
  \begin{tikzpicture}[remember picture,overlay,x=1bp,y=-1bp]
  \begin{scope}[shift={(current page.north west)}]
  \Fundo{19}
  \Texto{40.175}{28.518}{204.857}{17.983}{b}{2C2084}{Média, mediana e moda}
  \Texto{85.806}{83.374}{5.262}{9.465}{i}{262848}{n}
  \Texto{99.379}{64.087}{3.684}{6.625}{i}{262848}{n}
  \Texto{251.939}{87.561}{129.622}{6.974}{t}{000000}{x\ =\ np.array([2,\ 3,\ 4,\ 4,\ 5,\ 6,\ 8])}
  \Texto{251.939}{105.494}{29.625}{6.974}{t}{000000}{x.mean()}
  \Texto{337.13}{105.494}{29.625}{6.974}{c}{737696}{\#\ 4.5714}
  \Texto{251.939}{114.46}{44.437}{6.974}{t}{000000}{np.median(x)}
  \Texto{337.13}{114.46}{18.515}{6.974}{c}{737696}{\#\ 4.0}
  \Texto{251.939}{123.427}{70.359}{6.974}{t}{000000}{pd.Series(x).mode()}
  \Texto{337.13}{123.427}{11.109}{6.974}{c}{737696}{\#\ 4}
  \Texto{49.227}{168.935}{36.287}{8.518}{r}{000000}{regra que}
  \Texto{87.881}{168.935}{42.366}{8.966}{t}{2C2084}{np.median}
  \Texto{132.605}{168.935}{201.655}{8.518}{r}{000000}{aplica automaticamente. A moda pode não ser única.}
  \Texto{37.21}{247.463}{40.499}{6.625}{b}{737696}{EMASC 2026}
  \Texto{80.653}{247.463}{2.319}{6.625}{r}{FFB014}{\textbullet{}}
  \Texto{85.913}{247.463}{117.821}{6.625}{r}{737696}{Liga Acadêmica Matemática \& Empresa}
  \Texto{206.682}{247.463}{2.319}{6.625}{r}{FFB014}{\textbullet{}}
  \Texto{211.943}{247.463}{13.25}{6.625}{r}{737696}{UFS}
  \Texto{397.549}{248.132}{7.367}{6.625}{b}{2C2084}{14}
  \Texto{406.02}{248.132}{10.315}{6.625}{r}{737696}{/ 40}
  \end{scope}
  \end{tikzpicture}
\end{frame}


% Páginas 20 a 21 do PDF de Gabriel
\begin{frame}[plain,noframenumbering,label=original-20]
  \begin{tikzpicture}[remember picture,overlay,x=1bp,y=-1bp]
  \begin{scope}[shift={(current page.north west)}]
  \only<1>{\Fundo{20}}
  \only<2>{\Fundo{21}}
  \Texto{40.175}{29.916}{221.239}{17.983}{b}{2C2084}{Média e valores extremos}
  \Texto{35.433}{63.506}{295.754}{8.518}{i}{737696}{Uma medida só é confiável se soubermos como ela reage a valores extremos.}
  \only<1>{\Texto{49.227}{226.584}{7.104}{8.518}{r}{000000}{4.}}
  \Texto{37.21}{247.463}{40.499}{6.625}{b}{737696}{EMASC 2026}
  \Texto{80.653}{247.463}{2.319}{6.625}{r}{FFB014}{\textbullet{}}
  \Texto{85.913}{247.463}{117.821}{6.625}{r}{737696}{Liga Acadêmica Matemática \& Empresa}
  \Texto{206.682}{247.463}{2.319}{6.625}{r}{FFB014}{\textbullet{}}
  \Texto{211.943}{247.463}{13.25}{6.625}{r}{737696}{UFS}
  \Texto{397.549}{248.132}{7.367}{6.625}{b}{2C2084}{15}
  \Texto{406.02}{248.132}{10.315}{6.625}{r}{737696}{/ 40}
  \only<2>{\Texto{49.227}{226.584}{38.348}{8.518}{r}{000000}{A média é}}
  \only<2>{\Texto{89.942}{226.584}{33.621}{8.518}{b}{000000}{sensível}}
  \only<2>{\Texto{125.93}{226.584}{136.961}{8.518}{r}{000000}{a valores extremos; a mediana, não.}}
  \end{scope}
  \end{tikzpicture}
\end{frame}


% Página 22 do PDF de Gabriel
\begin{frame}[plain,noframenumbering,label=original-22]
  \begin{tikzpicture}[remember picture,overlay,x=1bp,y=-1bp]
  \begin{scope}[shift={(current page.north west)}]
  \Fundo{22}
  \Texto{40.175}{27.942}{124.906}{17.983}{b}{2C2084}{Desvio padrão}
  \Texto{35.433}{63.506}{264.723}{8.518}{i}{737696}{O valor típico não diz se os dados estão concentrados ou espalhados.}
  \Texto{49.227}{209.394}{357.083}{8.518}{r}{000000}{A soma dos desvios em relação à média é zero. Elevar ao quadrado elimina o sinal e a raiz}
  \Texto{37.21}{247.463}{40.499}{6.625}{b}{737696}{EMASC 2026}
  \Texto{80.653}{247.463}{2.319}{6.625}{r}{FFB014}{\textbullet{}}
  \Texto{85.913}{247.463}{117.821}{6.625}{r}{737696}{Liga Acadêmica Matemática \& Empresa}
  \Texto{206.682}{247.463}{2.319}{6.625}{r}{FFB014}{\textbullet{}}
  \Texto{211.943}{247.463}{13.25}{6.625}{r}{737696}{UFS}
  \Texto{397.549}{248.132}{7.367}{6.625}{b}{2C2084}{16}
  \Texto{406.02}{248.132}{10.315}{6.625}{r}{737696}{/ 40}
  \end{scope}
  \end{tikzpicture}
\end{frame}


% Página 23 do PDF de Gabriel
\begin{frame}[plain,noframenumbering,label=original-23]
  \begin{tikzpicture}[remember picture,overlay,x=1bp,y=-1bp]
  \begin{scope}[shift={(current page.north west)}]
  \Fundo{23}
  \Texto{40.175}{28.167}{223.792}{17.983}{b}{2C2084}{Variância e desvio padrão}
  \Texto{102.646}{69.8}{5.262}{9.465}{r}{262848}{1}
  \Texto{124.86}{64.087}{3.684}{6.625}{i}{262848}{n}
  \Texto{113.253}{116.688}{5.262}{9.465}{r}{262848}{6}
  \Texto{251.939}{87.561}{59.256}{6.974}{t}{000000}{d\ =\ x\ -\ x.mean()}
  \Texto{251.939}{96.528}{25.922}{6.974}{t}{000000}{(d**2).}
  \Texto{277.861}{96.528}{11.109}{6.974}{t}{2C2084}{sum}
  \Texto{288.97}{96.528}{7.406}{6.974}{t}{000000}{()}
  \Texto{318.61}{96.528}{29.625}{6.974}{c}{737696}{\#\ 23.714}
  \Texto{251.939}{105.494}{48.14}{6.974}{t}{000000}{x.var(ddof=1)}
  \Texto{318.61}{105.494}{25.922}{6.974}{c}{737696}{\#\ 3.952}
  \Texto{251.939}{114.46}{48.14}{6.974}{t}{000000}{x.std(ddof=1)}
  \Texto{318.61}{114.46}{25.922}{6.974}{c}{737696}{\#\ 1.988}
  \Texto{49.227}{153.918}{44.49}{8.518}{r}{000000}{Em NumPy,}
  \Texto{96.008}{153.918}{28.244}{8.966}{t}{2C2084}{ddof=0}
  \Texto{126.518}{153.918}{211.664}{8.518}{r}{000000}{calcula a variância populacional; em pandas, o padrão é}
  \Texto{340.456}{153.918}{28.244}{8.966}{t}{2C2084}{ddof=1}
  \Texto{368.694}{153.918}{37.616}{8.518}{r}{000000}{, a variân-}
  \Texto{49.227}{164.877}{319.297}{8.518}{r}{000000}{cia amostral. Para estimar a variância da população a partir de uma amostra, usa-se}
  \Texto{370.888}{164.877}{28.244}{8.966}{t}{2C2084}{ddof=1}
  \Texto{399.125}{164.877}{2.368}{8.518}{r}{000000}{.}
  \Texto{37.21}{247.463}{40.499}{6.625}{b}{737696}{EMASC 2026}
  \Texto{80.653}{247.463}{2.319}{6.625}{r}{FFB014}{\textbullet{}}
  \Texto{85.913}{247.463}{117.821}{6.625}{r}{737696}{Liga Acadêmica Matemática \& Empresa}
  \Texto{206.682}{247.463}{2.319}{6.625}{r}{FFB014}{\textbullet{}}
  \Texto{211.943}{247.463}{13.25}{6.625}{r}{737696}{UFS}
  \Texto{397.549}{248.132}{7.367}{6.625}{b}{2C2084}{17}
  \Texto{406.02}{248.132}{10.315}{6.625}{r}{737696}{/ 40}
  \end{scope}
  \end{tikzpicture}
\end{frame}


% Página 24 do PDF de Gabriel
\begin{frame}[plain,noframenumbering,label=original-24]
  \begin{tikzpicture}[remember picture,overlay,x=1bp,y=-1bp]
  \begin{scope}[shift={(current page.north west)}]
  \Fundo{24}
  \Texto{35.433}{63.506}{282.891}{8.518}{i}{737696}{Estimar a variância da população a partir de uma amostra exige um ajuste.}
  \Texto{49.227}{181.15}{296.29}{8.518}{r}{000000}{A média é calculada da própria amostra, o que reduz os desvios. Dividir por}
  \Texto{348.536}{181.15}{4.736}{8.518}{i}{000000}{n}
  \Texto{356.428}{181.15}{42.13}{8.518}{b}{000000}{subestima}
  \Texto{401.573}{181.15}{4.736}{8.518}{r}{000000}{a}
  \Texto{37.21}{247.463}{40.499}{6.625}{b}{737696}{EMASC 2026}
  \Texto{80.653}{247.463}{2.319}{6.625}{r}{FFB014}{\textbullet{}}
  \Texto{85.913}{247.463}{117.821}{6.625}{r}{737696}{Liga Acadêmica Matemática \& Empresa}
  \Texto{206.682}{247.463}{2.319}{6.625}{r}{FFB014}{\textbullet{}}
  \Texto{211.943}{247.463}{13.25}{6.625}{r}{737696}{UFS}
  \Texto{397.549}{248.132}{7.367}{6.625}{b}{2C2084}{18}
  \Texto{406.02}{248.132}{10.315}{6.625}{r}{737696}{/ 40}
  \end{scope}
  \end{tikzpicture}
\end{frame}


% Página 25 do PDF de Gabriel
\begin{frame}[plain,noframenumbering,label=original-25]
  \begin{tikzpicture}[remember picture,overlay,x=1bp,y=-1bp]
  \begin{scope}[shift={(current page.north west)}]
  \Fundo{25}
  \Texto{40.175}{29.336}{117.282}{17.983}{b}{2C2084}{Quartis e IQR}
  \Texto{35.433}{63.506}{298.139}{8.518}{i}{737696}{Quando há valores extremos, é útil medir a dispersão sem depender da média.}
  \Texto{49.227}{142.749}{130.479}{8.518}{r}{000000}{respondem a 25\%, 50\% e 75\%. O}
  \Texto{182.209}{142.749}{15.145}{8.518}{b}{000000}{IQR}
  \Texto{199.859}{142.749}{138.179}{8.518}{r}{000000}{é a amplitude dos 50\% centrais e as}
  \Texto{340.541}{142.749}{26.823}{8.518}{b}{000000}{cercas}
  \Texto{369.87}{142.749}{36.44}{8.518}{r}{000000}{delimitam}
  \Texto{37.21}{247.463}{40.499}{6.625}{b}{737696}{EMASC 2026}
  \Texto{80.653}{247.463}{2.319}{6.625}{r}{FFB014}{\textbullet{}}
  \Texto{85.913}{247.463}{117.821}{6.625}{r}{737696}{Liga Acadêmica Matemática \& Empresa}
  \Texto{206.682}{247.463}{2.319}{6.625}{r}{FFB014}{\textbullet{}}
  \Texto{211.943}{247.463}{13.25}{6.625}{r}{737696}{UFS}
  \Texto{397.549}{248.132}{7.367}{6.625}{b}{2C2084}{19}
  \Texto{406.02}{248.132}{10.315}{6.625}{r}{737696}{/ 40}
  \end{scope}
  \end{tikzpicture}
\end{frame}
